\section*{Capítulo \ref{ch:g6:exploring-our-environment} --- Explorar o nosso ambiente}

\begin{solution}{exo:g6:exploring-our-environment:1}
\omterm{def:g6:exploring-our-environment:components}{Seres vivos} (musgo, tatuzinhos); \omterm{def:g6:exploring-our-environment:components}{matéria mineral} (a pedra do muro, a
poça); vestígios e restos de \omterm{def:g6:exploring-our-environment:components}{seres vivos} (\omterm{def:g1:plants-around-us:parts}{folhas} secas, a \omterm{def:g1:animals-around-us:covering}{concha} de
caracol).
\end{solution}

\begin{solution}{exo:g6:exploring-our-environment:2}
\omterm{def:g6:exploring-our-environment:conditions}{Temperatura}, com um termômetro (em graus Celsius); \omterm{def:g6:exploring-our-environment:conditions}{umidade} do ar e
do solo; \omterm{def:g6:exploring-our-environment:conditions}{luz}, do sol pleno à sombra profunda.
\end{solution}

\begin{solution}{exo:g6:exploring-our-environment:3}
Vivos: a samambaia, a centopeia. Mineral: a poça, a pedra do muro.
Vestígios: a \omterm{def:g1:animals-around-us:covering}{concha} de caracol, as fezes de ave.
\end{solution}

\begin{solution}{exo:g6:exploring-our-environment:4}
Dezesseis graus (\qty{38}{\degreeCelsius} contra
\qty{22}{\degreeCelsius}) e sol pleno contra sombra profunda (seco
contra úmido, também). O matinho da seca e do brilho forte combina
com o pé do muro ensolarado; musgo, \omterm{prop:g4:plants-without-seeds:damp}{samambaias} e tatuzinhos combinam
com o canto fresco e úmido.
\end{solution}

\begin{solution}{exo:g6:exploring-our-environment:5}
O revestimento dele deixa a água escapar, então, no sol seco, o
corpo dele perde a água depressa demais e isso é fatal. É por isso
que os tatuzinhos são encontrados em \omterm{def:g6:exploring-our-environment:microhabitat}{micro-hábitats} úmidos e escuros
--- debaixo de pedras, debaixo de madeira caída.
\end{solution}

\begin{solution}{exo:g6:exploring-our-environment:6}
Por exemplo: uma avelã roída --- o rato-do-mato; uma teia --- a
aranha dela; pegadas na lama --- a garça (fezes, \omterm{def:g1:animals-around-us:covering}{penas} e \omterm{def:g1:animals-around-us:covering}{conchas}
servem do mesmo jeito).
\end{solution}

\begin{solution}{exo:g6:exploring-our-environment:7}
Uma caixa fechada, metade com papel úmido e metade com papel seco,
tatuzinhos soltos na linha do meio; minutos depois quase todos estão
do lado úmido (e o mesmo vale para sombra contra \omterm{def:g6:exploring-our-environment:conditions}{luz}). Isso mostra
que cada \omterm{def:g1:animals-around-us:animal}{animal} se instala onde as \omterm{def:g6:exploring-our-environment:conditions}{condições} combinam com o corpo
dele --- o mapa do grupo se desenha sozinho, sem ninguém conduzir.
\end{solution}

\begin{solution}{exo:g6:exploring-our-environment:8}
Porque só medidas feitas do mesmo jeito podem ser comparadas: uma
diferença precisa vir dos sítios, e não da medição. É a versão de
campo do teste justo --- uma coisa varia (o sítio) e todo o resto é
mantido igual.
\end{solution}

\begin{solution}{exo:g6:exploring-our-environment:9}
Topo da pedra: iluminado, seco, quente --- líquen (com a mosca
tomando sol como visitante). Parte de baixo da pedra: escura, úmida,
fresca --- tatuzinhos (ou a centopeia). A cinco centímetros de
distância, dois conjuntos combinados de moradores.
\end{solution}

\begin{solution}{exo:g6:exploring-our-environment:10}
Passo 1: marcar um metro quadrado definido. Passo 2: medir --- a
leitura do termômetro anotada, a \omterm{def:g6:exploring-our-environment:conditions}{umidade} do solo sentida e anotada,
a classe de \omterm{def:g6:exploring-our-environment:conditions}{luz} anotada. Passo 3: vasculhar direito, com as pedras
levantadas e recolocadas, listando os três tipos de componentes.
Passo 4: registrar ali mesmo --- contagens, e não ``uns''. Passo 5:
com um sítio só, nenhuma comparação era sequer possível; levantar um
segundo sítio do mesmo jeito.
\end{solution}

\begin{solution}{exo:g6:exploring-our-environment:11}
Explicação: a hera prospera na faixa de \omterm{def:g6:exploring-our-environment:conditions}{condições} do muro sombreado
(sombra, \omterm{def:g6:exploring-our-environment:conditions}{umidade}) e não na do muro ensolarado (brilho forte, seca).
Teste: medir o pé dos dois muros --- \omterm{def:g6:exploring-our-environment:conditions}{temperatura}, \omterm{def:g6:exploring-our-environment:conditions}{umidade}, \omterm{def:g6:exploring-our-environment:conditions}{luz} --- e
conferir com as preferências conhecidas da hera; ou procurar hera em
outros muros e registrar que orientação esses muros têm em comum.
\end{solution}

\begin{solution}{exo:g6:exploring-our-environment:12}
Elimine a diferença de parede: faça a escolha numa caixa redonda
(parede igual em toda parte), com as metades úmida e seca trocadas
entre uma rodada e outra --- úmida a leste na rodada 1, úmida a oeste
na rodada 2. Se os tatuzinhos seguirem a \omterm{def:g6:exploring-our-environment:conditions}{umidade} para onde quer que
ela seja posta, as paredes estão descartadas; entre as rodadas só a
condição testada mudou, e todo o resto foi mantido igual.
\end{solution}

\begin{solution}{pb:g6:exploring-our-environment:1}
\textbf{1.} A cerca viva alta: ela dá sombra à vala N (e a abriga),
enquanto a vala S fica virada para o sol o dia inteiro.

\textbf{2.} \omterm{def:g6:exploring-our-environment:conditions}{Temperatura}: \qty{19}{\degreeCelsius} contra
\qty{31}{\degreeCelsius}. \omterm{def:g6:exploring-our-environment:conditions}{Umidade}: a terra da vala N está úmida
(gruda), a da vala S está seca (vira pó). A \omterm{def:g6:exploring-our-environment:conditions}{luz} completa o trio:
sombra contra sol pleno.

\textbf{3.} As \omterm{def:g6:exploring-our-environment:conditions}{condições} mudam ao longo do dia: às 18 horas qualquer
sítio está mais fresco e com menos \omterm{def:g6:exploring-our-environment:conditions}{luz}, então a diferença medida
misturaria sítio com hora. A mesma hora mantém a comparação justa.

\textbf{4.} Para que o esforço de busca seja igual: mais área ou mais
tempo num dos sítios encontraria mais moradores ali --- um teste
injusto de qual das valas é mais rica.

\textbf{5.} As duas listas nomeiam só \omterm{def:g6:exploring-our-environment:components}{seres vivos}. Faltam: os
componentes minerais (solo, pedras, água) e os vestígios --- que um
inventário direito também precisa registrar.

\textbf{6.} Tatuzinhos: vala N --- eles precisam de \omterm{def:g6:exploring-our-environment:conditions}{umidade} e sombra,
senão ressecam. O sapo: vala N --- \omterm{def:g1:the-five-senses:sense}{pele} úmida, precisando de abrigo
úmido. O lagarto: vala S --- ele esquenta o corpo no sol e caça onde
pode se aquecer.

\textbf{7.} Ou existem caracóis morando ali escondidos (talvez ativos
só nas horas úmidas da noite), ou a \omterm{def:g1:animals-around-us:covering}{concha} é um resto vindo de outro
lugar ou de outra estação --- um vestígio prova presença \emph{em
algum momento}, e não moradia agora.

\textbf{8.} Escuro, úmido, fresco. A regra da
\cref{prop:g6:exploring-our-environment:decide}: cada ponto é
habitado pelas \omterm{def:g5:biodiversity:species}{espécies} cujas faixas combinam com as \omterm{def:g6:exploring-our-environment:conditions}{condições} dele.

\textbf{9.} A N tem musgo e sapos porque a sombra a mantém fresca e
úmida, coisa que o corpo deles exige; a S tem lagartos e matagais de
\omterm{def:g1:plants-around-us:parts}{folhas} enceradas porque o sol pleno a deixa quente e seca, o que
serve a um caçador que toma sol e a uma \omterm{def:g1:plants-around-us:plant}{planta} à prova de seca.

\textbf{10.} Sem a cerca viva, a vala N recebe sol pleno: mais
quente, mais seca --- \omterm{def:g6:exploring-our-environment:conditions}{condições} caminhando para as da vala S de hoje.
Espere que os amantes da \omterm{def:g6:exploring-our-environment:conditions}{umidade} (musgo, \omterm{prop:g4:plants-without-seeds:damp}{samambaias}, tatuzinhos, o
sapo) diminuam e que os amantes do sol (capins secos, gafanhotos e
talvez o lagarto) se mudem para lá.

\textbf{11.} Porque o novo levantamento precisa diferir do primeiro
em uma coisa só --- a ausência da cerca viva. Sítios, instrumentos ou
estação novos acrescentariam diferenças próprias, e a mudança já não
poderia ser atribuída à cerca viva: é a regra do teste justo na
escala de uma paisagem.

\textbf{12.} Cada tatuzinho anda e vira até o corpo dele ficar
confortável, então o grupo acaba reunido onde está a \omterm{def:g6:exploring-our-environment:conditions}{umidade} --- as
posições dos \omterm{def:g1:animals-around-us:animal}{animais} registram as \omterm{def:g6:exploring-our-environment:conditions}{condições}, com a mesma segurança
que as leituras de um termômetro. Ninguém os conduz; o conforto de
cada indivíduo é que desenha o mapa, e é justamente por isso que os
moradores podem ser lidos como medidas.
\end{solution}
